\documentclass[11pt,a4paper]{article}
\usepackage{fontspec,amsmath,amssymb,amsthm,microtype}
\setmainfont{lmroman10-regular.otf}[BoldFont=lmroman10-bold.otf,ItalicFont=lmroman10-italic.otf,BoldItalicFont=lmroman10-bolditalic.otf]
\usepackage{xcolor}
\usepackage[most]{tcolorbox}
\usepackage{mathrsfs}
\providecommand{\lt}{<}
\providecommand{\gt}{>}
\setmonofont{DejaVuSansMono.ttf}[BoldFont=DejaVuSansMono-Bold.ttf,ItalicFont=DejaVuSansMono-Oblique.ttf,BoldItalicFont=DejaVuSansMono-BoldOblique.ttf,Scale=0.86]
\newfontfamily\leanfallback{STIXTwoMath-Regular.otf}[Scale=0.86]
\definecolor{leanrule}{gray}{0.62}
\definecolor{leanback}{gray}{0.962}
\definecolor{leantitle}{gray}{0.22}
\definecolor{leanmuted}{gray}{0.38}
\newtcolorbox{leanbox}[1]{enhanced,breakable,colback=leanback,colframe=leanrule,
  boxrule=0.4pt,arc=0.6mm,left=2.4mm,right=2.4mm,top=1.8mm,bottom=1.8mm,
  before skip=0.8em,after skip=1em,fonttitle=\small\bfseries,coltitle=leantitle,
  colbacktitle=leanback,titlerule=0.3pt,title={#1}}
\newlength{\leancharwidth}
\newenvironment{leancode}{\par\smallskip\begingroup\ttfamily\footnotesize\raggedright
  \settowidth{\leancharwidth}{M}\setlength{\parskip}{0pt}\setlength{\parindent}{0pt}
  \frenchspacing}{\par\endgroup\smallskip}
\newcommand{\leanline}[2]{\hangindent=\dimexpr#1\leancharwidth+4\leancharwidth\relax
  \hangafter=1\noindent\hspace*{#1\leancharwidth}#2\par}
\newcommand{\leansource}[1]{\par{\footnotesize\color{leanmuted}#1}\par}
\newcommand{\leanlabel}[1]{\par\smallskip\noindent{\small\bfseries #1}\par}
\usepackage[margin=25mm]{geometry}
\usepackage[unicode,hidelinks]{hyperref}
\hypersetup{pdfauthor={LangevinLean Project}}
\setlength{\emergencystretch}{3em}
\setlength{\parskip}{0.4em}
\setlength{\parindent}{0pt}
\allowdisplaybreaks[2]
\raggedbottom
\title{\texorpdfstring{The Ornstein–Uhlenbeck innovation law and independence across time}{The Ornstein–Uhlenbeck innovation law and independence across time}}
\author{LangevinLean Project}
\date{}
\begin{document}\maketitle
Draft mathematical exposition; these notes may contain mistakes.


The constructed convolution is a random variable on the constructed Brownian source. This section proves that its fluctuation-dissipation scaling has exactly thermal Gaussian law, and that full spatial vectors on distinct grid intervals are independent. The complete stream therefore has the independent Gaussian product law required by the Markov transition.



The thermal response law is computed from its exponential kernel and assembled across independent spatial coordinates. Disjoint intervals yield the complete independent Gaussian stream required by the splitting chain.


\section{\texorpdfstring{Identify independent spatial copies through their scalar law}{Identify independent spatial copies through their scalar law}}


A scalar \(L^{2}\) variable \(Z\) is evaluated on each independent coordinate Brownian source. Each copy has the same proved law \(\nu\), and the coordinate product construction proves their independence. The infinite-product law theorem, specialized to finitely many spatial coordinates, gives the finite product measure \(\nu^{\otimes\iota}\). Transporting this law through the Euclidean coordinate map identifies the law of the coordinate assembly map; its representative theorem supplies the almost-everywhere equality. For \(\nu\) equal to a scaled real standard Gaussian, the product pushforward identity and the definition-independent standard Gaussian theorem identify this law as the same scalar scaling of the Euclidean standard Gaussian.



The scalar response at each spatial index has the computed centered Gaussian law. Independence of the coordinate sources gives their product law. Assembling coordinates then yields a Euclidean Gaussian with a scalar covariance multiple of the identity.

\begin{leanbox}{Lean: \texttt{spatialLift\_\allowbreak{}hasLaw}, \texttt{spatialLift\_\allowbreak{}hasLaw\_\allowbreak{}scaledStandard}}
\leansource{\texttt{spatialLift\_\allowbreak{}hasLaw} \quad KineticLangevin/\allowbreak{}OUMarginal.lean, lines 13–30}
\begin{leancode}
\leanline{0}{theorem\ spatialLift\_hasLaw\ (Z\ :\ Lp\ ℝ\ 2\ probability)\ (ν\ :\ Measure\ ℝ)}
\leanline{4}{[IsProbabilityMeasure\ ν]\ (hZ\ :\ HasLaw\ Z\ ν\ probability)\ :}
\leanline{4}{HasLaw\ (spatialLift\ ι\ Z)\ ((Measure.pi\ (fun\ \_\ :\ ι\ ↦\ ν)).map\ (WithLp.toLp\ 2))}
\leanline{6}{(ULA.CanonicalBrownian.vectorProbability\ ι)}
\end{leancode}
\leansource{\texttt{spatialLift\_\allowbreak{}hasLaw\_\allowbreak{}scaledStandard} \quad KineticLangevin/\allowbreak{}OUMarginal.lean, lines 34–46}
\begin{leancode}
\leanline{0}{theorem\ spatialLift\_hasLaw\_scaledStandard\ (Z\ :\ Lp\ ℝ\ 2\ probability)\ (σ\ :\ ℝ)}
\leanline{4}{(hZ\ :\ HasLaw\ Z\ ((gaussianReal\ 0\ 1).map\ (fun\ x\ ↦\ σ*x))\ probability)\ :}
\leanline{4}{HasLaw\ (spatialLift\ ι\ Z)}
\leanline{6}{((stdGaussian\ (EuclideanSpace\ ℝ\ ι)).map\ (fun\ x\ ↦\ σ\ •\ x))}
\leanline{6}{(ULA.CanonicalBrownian.vectorProbability\ ι)}
\end{leancode}
\leanlabel{How the Lean proof follows the argument}
\texttt{spatial\allowbreak{}Lift\_\allowbreak{}has\allowbreak{}Law} gives the formal counterpart of this step. The main intermediate facts are \texttt{hm}, \texttt{hi}, \texttt{hl}, \texttt{hp}, \texttt{ht} in that order.




\texttt{spatial\allowbreak{}Lift\_\allowbreak{}has\allowbreak{}Law\_\allowbreak{}scaled\allowbreak{}Standard} gives the formal counterpart of this step. The main intermediate facts are \texttt{hprod} in that order. Named proof applications include \texttt{spatial\allowbreak{}Lift\_\allowbreak{}has\allowbreak{}Law}; the adjacent proof shows their hypotheses and order of use.
\end{leanbox}


Earlier results used here: \cite[Section~3.8.16]{LangevinLean2026ULASC}; \cite[Section~3.8.17]{LangevinLean2026ULASC}.

\section{\texorpdfstring{Compute the scalar thermal variance}{Compute the scalar thermal variance}}


The scalar segment convolution is the constructed isonormal map applied to its time kernel \(K\), so it has centered Gaussian variance \(\|K\|^2\). With \(c=\sqrt{2\gamma}\), its scaled variance is \(2\gamma\|K\|^2\). The proved kernel inner-product identity gives \(\|K\|^2=\int_0^h e^{-2\gamma s}\,ds\), and the exact fluctuation-decay balance gives \(2\gamma\|K\|^2=1-e^{-2\gamma h}\). This equals \(\sigma_h^2\) for \(\sigma_h=\sqrt{1-e^{-2\gamma h}}\). The argument derives nonnegativity from the damping bound, uses the isonormal Gaussian law, checks the equality of nonnegative-real variances, and rewrites the scaled real Gaussian pushforward. Zero friction is included without division.



The filtered velocity response variance is the square integral of the exponential kernel. After multiplying by \(2\gamma\), it is \(1-e^{-2\gamma h}\). This is the missing variance of a velocity damped for time \(h\).

\begin{leanbox}{Lean: \texttt{scalarOU\_\allowbreak{}hasLaw}}
\leansource{\texttt{scalarOU\_\allowbreak{}hasLaw} \quad KineticLangevin/\allowbreak{}OUMarginal.lean, lines 50–80}
\begin{leancode}
\leanline{0}{theorem\ scalarOU\_hasLaw\ (γ\ :\ ℝ)\ (hγ\ :\ 0≤γ)\ (a\ t\ :\ ℝ≥0)\ :}
\leanline{4}{HasLaw\ (Real.sqrt\ (2*γ)\ •\ segmentConvolutionLp\ (psi\ γ\ 0)\ (continuous\_psi\ γ\ 0)\ a\ t)}
\leanline{6}{((gaussianReal\ 0\ 1).map}
\leanline{8}{(fun\ x\ ↦\ Real.sqrt\ (1-(psi\ γ\ 0\ t)\textasciicircum{}2)*x))\ probability}
\end{leancode}
\leanlabel{How the Lean proof follows the argument}
\texttt{scalar\allowbreak{}OU\_\allowbreak{}has\allowbreak{}Law} gives the formal counterpart of this step. The main intermediate facts are \texttt{K}, \texttt{c}, \texttt{\ensuremath{\sigma}}, \texttt{hvariance}, \texttt{hnorm}, \texttt{hnormc} in that order. Named proof applications include \texttt{NNReal.\allowbreak{}eq}, \texttt{Real.\allowbreak{}coe\_\allowbreak{}to\allowbreak{}NNReal}, \texttt{hz}; the adjacent proof shows their hypotheses and order of use.
\end{leanbox}


Earlier results used here: \cite[Section~3.6.12]{LangevinLean2026ULASC}; \cite[Section~3.6.24]{LangevinLean2026ULASC}; \cite[Section~3.8.5]{LangevinLean2026ULASC}.

\section{\texorpdfstring{The entire spatial OU Gaussian law}{The entire spatial OU Gaussian law}}


Apply the independent-spatial-copy result to the proved scalar thermal law. Linearity of the coordinate assembly map identifies its scaled scalar input with the scaled vector segment convolution. Thus \(\sqrt{2\gamma}\,Z_{a,h}\) has exactly the law of \(\sqrt{1-e^{-2\gamma h}}\,\xi\), where \(\xi\) is the Euclidean standard Gaussian. This proves the full vector distribution rather than only its covariance.



A centered Euclidean Gaussian is determined by its covariance bilinear form. The independent coordinates give covariance equal to the scalar thermal variance times the identity. Thus the assembled innovation is exactly the thermal Gaussian law required by the thermostat.

\begin{leanbox}{Lean: \texttt{vectorOU\_\allowbreak{}hasLaw}}
\leansource{\texttt{vectorOU\_\allowbreak{}hasLaw} \quad KineticLangevin/\allowbreak{}OUMarginal.lean, lines 84–92}
\begin{leancode}
\leanline{0}{theorem\ vectorOU\_hasLaw\ (γ\ :\ ℝ)\ (hγ\ :\ 0≤γ)\ (a\ t\ :\ ℝ≥0)\ :}
\leanline{4}{HasLaw\ (Real.sqrt\ (2*γ)\ •\ vectorSegmentLp\ ι\ (psi\ γ\ 0)\ (continuous\_psi\ γ\ 0)\ a\ t)}
\leanline{6}{((stdGaussian\ (EuclideanSpace\ ℝ\ ι)).map}
\leanline{8}{(fun\ x\ ↦\ Real.sqrt\ (1-(psi\ γ\ 0\ t)\textasciicircum{}2)\ •\ x))}
\leanline{6}{(ULA.CanonicalBrownian.vectorProbability\ ι)}
\end{leancode}
\leanlabel{How the Lean proof follows the argument}
\texttt{vector\allowbreak{}OU\_\allowbreak{}has\allowbreak{}Law} gives the formal counterpart of this step. The main intermediate facts are \texttt{hz} in that order. Named proof applications include \texttt{hz}; the adjacent proof shows their hypotheses and order of use.
\end{leanbox}


Earlier results used here: \cite[Section~3.8.17]{LangevinLean2026ULASC}.

\section{\texorpdfstring{Independence of full-vector innovations across time}{Independence of full-vector innovations across time}}


The constructed KLMC innovation pairs are already proved independent over disjoint time-grid intervals by their Brownian kernel supports. Project each pair to its velocity component and apply any fixed scalar scaling. The \(L^{2}\) vector-segment representatives agree almost everywhere with these projected constructed innovations. Independence is preserved under these measurable projections and almost-everywhere changes. The argument explicitly combines the representative scalar-multiplication event with the constructed terminal state identity.



Entire vector blocks from different intervals are measurable functions of disjoint collections of Gaussian response variables. Independence of those collections therefore gives independence of the vectors, rather than only pairwise independence of their individual coordinates.

\begin{leanbox}{Lean: \texttt{vectorOU\_\allowbreak{}grid\_\allowbreak{}independent}}
\leansource{\texttt{vectorOU\_\allowbreak{}grid\_\allowbreak{}independent} \quad KineticLangevin/\allowbreak{}OUMarginal.lean, lines 96–107}
\begin{leancode}
\leanline{0}{theorem\ vectorOU\_grid\_independent\ (d\ :\ ℕ)\ (γ\ c\ :\ ℝ)\ (h\ :\ ℝ≥0)\ :}
\leanline{4}{iIndepFun\ (fun\ n\ :\ ℕ\ ↦\ c\ •\ vectorSegmentLp\ (Fin\ d)\ (psi\ γ\ 0)\ (continuous\_psi\ γ\ 0)}
\leanline{6}{((n\ :\ ℝ≥0)*h)\ h)\ (ULA.CanonicalBrownian.vectorProbability\ (Fin\ d))}
\end{leancode}
\leanlabel{How the Lean proof follows the argument}
\texttt{vector\allowbreak{}OU\_\allowbreak{}grid\_\allowbreak{}independent} gives the formal counterpart of this step. The main intermediate facts are \texttt{hi} in that order. Named proof applications include \texttt{hi.\allowbreak{}congr}; the adjacent proof shows their hypotheses and order of use.
\end{leanbox}


Earlier results used here: \cite[Section~3.8.17]{LangevinLean2026ULASC}.

\section{\texorpdfstring{The full independent Gaussian stream}{The full independent Gaussian stream}}


Every scaled OU grid innovation has the exact thermal Gaussian law, and the preceding theorem proves independence of all grid vectors. The infinite-product law theorem therefore identifies the entire innovation stream with the product of this Gaussian law over all iterations. Measurability follows from the strongly measurable \(L^{2}\) representative at each countable grid index. This stream identity is the law identity used to transfer the Brownian-driven splitting chain to its iterated Markov laws.



The finite-dimensional distributions of the full stream are products of thermal Gaussian vector law. This identifies the stream as independent and identically distributed. It supplies the exact fresh-noise law required when the splitting transition is iterated.

\begin{leanbox}{Lean: \texttt{vectorOU\_\allowbreak{}grid\_\allowbreak{}hasLaw}}
\leansource{\texttt{vectorOU\_\allowbreak{}grid\_\allowbreak{}hasLaw} \quad KineticLangevin/\allowbreak{}OUMarginal.lean, lines 111–124}
\begin{leancode}
\leanline{0}{theorem\ vectorOU\_grid\_hasLaw\ (d\ :\ ℕ)\ (γ\ :\ ℝ)\ (hγ\ :\ 0≤γ)\ (h\ :\ ℝ≥0)\ :}
\leanline{4}{HasLaw\ (fun\ ω\ (n\ :\ ℕ)\ ↦\ (Real.sqrt\ (2*γ)\ •}
\leanline{6}{vectorSegmentLp\ (Fin\ d)\ (psi\ γ\ 0)\ (continuous\_psi\ γ\ 0)\ ((n\ :\ ℝ≥0)*h)\ h)\ ω)}
\leanline{6}{(Measure.infinitePi\ (fun\ \_\ :\ ℕ\ ↦}
\leanline{8}{(stdGaussian\ (EuclideanSpace\ ℝ\ (Fin\ d))).map}
\leanline{10}{(fun\ x\ ↦\ Real.sqrt\ (1-(psi\ γ\ 0\ h)\textasciicircum{}2)\ •\ x)))}
\leanline{6}{(ULA.CanonicalBrownian.vectorProbability\ (Fin\ d))}
\end{leancode}
\leanlabel{How the Lean proof follows the argument}
\texttt{vector\allowbreak{}OU\_\allowbreak{}grid\_\allowbreak{}has\allowbreak{}Law} gives the formal counterpart of this step. Named proof applications include \texttt{vector\allowbreak{}OU\_\allowbreak{}has\allowbreak{}Law}, \texttt{measurable\_\allowbreak{}pi\_\allowbreak{}iff.\allowbreak{}mpr}; the adjacent proof shows their hypotheses and order of use.
\end{leanbox}


Earlier results used here: \cite[Section~3.8.17]{LangevinLean2026ULASC}.

\begin{thebibliography}{ULA-LSI}
\bibitem[ULA-SC]{LangevinLean2026ULASC}
LangevinLean Project. \emph{Unadjusted Langevin under Strong Convexity}. LL-01, 2026. Draft report, content revision d50c0f64815b. \url{https://odakyildiz.github.io/LangevinLean/docs/pdf/UnadjustedLangevin/FullProof.pdf}.
\end{thebibliography}

\end{document}
